\section{Lean formalization architecture}

The archive contains two Lean blueprints.  They are deliberately separated
into a semantic algorithm layer and a deeply embedded object-theory layer.

\subsection{Layer 1: executable semantic algorithms}

The file \code{lean/FiniteContractSemantic.lean} defines:
\begin{itemize}[leftmargin=1.6em]
\item bounded Boolean folds and first-witness search;
\item finite trace execution;
\item reverse demand propagation on a topologically ordered finite graph;
\item fuelled backward viability and coherent finite paths;
\item bounded enumeration;
\item rational cubic bisection specifications.
\end{itemize}
Its intended theorems are ordinary Lean statements about \code{Nat},
\code{List}, and \code{Bool}.  They serve three purposes: test the definitions,
clarify invariants, and provide external semantics lemmas.  They are not the
iPRA proof.

\subsection{Layer 2: actual ProveIt interfaces}

The object blueprint imports:
\begin{lstlisting}[language={}]
import ConciseIPRA.PR
import ConciseIPRA.NativeBiInterpretation
\end{lstlisting}
and works in \code{LeanProofs.ConciseIPRA}.  The relevant existing declarations
are:
\begin{lstlisting}[language={}]
PR.zero   PR.succ   PR.proj   PR.comp   PR.prec
PTm.var   PTm.app
PRAAx.proj  PRAAx.comp  PRAAx.precZero  PRAAx.precSucc
PRAAx.succNeZero  PRAAx.succInj  PRAAx.ind
PRA.proj  PRA.comp  PRA.precZero  PRA.precSucc
PRA.qfInd
pra_iff_native
\end{lstlisting}

The first implementation milestone is a small object-level arithmetic
library: numerals, addition, multiplication, truncated subtraction, order
Booleans, pair/list coding, Boolean folds, and table lookup.  Each function is
a concrete \code{PR} term, and each defining equation is a
\(\Prf\,\PRAAx\) theorem composed from \code{PRA.comp},
\code{PRA.precZero}, and \code{PRA.precSucc}.

\subsection{Object trace structure}

A practical record is:
\begin{lstlisting}[language={}]
structure PRTraceSystem (p : Nat) where
  init    : PR p
  step    : PR (p + 2)   -- stage, state, parameters
  guard   : PR (p + 1)   -- state, parameters -> Boolean
  finish  : PR (p + 1)
  horizon : PR p
\end{lstlisting}
The exact arities can be adjusted to the repository's convention that the
recursion argument comes first.  The invariant is not a Lean predicate on
natural numbers; it is an \code{Fm PTm} whose atomic components assert that
the Boolean \code{guard} term equals \(1\).

The core target has the schematic shape:
\begin{lstlisting}[language={}]
theorem PRA.traceReflection
    (S : PRTraceSystem p)
    (hQF : (traceInv S).QF)
    (hBase : Gamma |-[(PRAAx)] baseFormula S)
    (hStep : stepContext S Gamma |-[(PRAAx)] stepFormula S)
    (hFinish : finishContext S Gamma |-[(PRAAx)] outputFormula S) :
    Gamma |-[(PRAAx)] finalFormula S := by
  -- instantiate PRA.qfInd with the trace invariant;
  -- use S.init/S.step defining equations;
  -- compose hBase, hStep, hFinish.
\end{lstlisting}
The actual file uses ASCII comments rather than claiming the notation above
already elaborates.

\subsection{Compilation milestones}

\begin{enumerate}[label=\textbf{M\arabic*.},leftmargin=3em]
\item \textbf{Boolean folds.}
      Define \code{prAllBelow}, \code{prAnyBelow}, and \code{prFirst}.
      Prove their zero/successor equations and object-level correctness.
\item \textbf{Coded tables.}
      Define pair, projections, list length, lookup, append, and update as
      \code{PR} terms.  Prove the lookup-after-update equations needed by the
      trace invariant.
\item \textbf{Generic trace.}
      Encode \code{run} by primitive recursion.  Prove its zero and successor
      equations using \code{PRA.precZero}/\code{PRA.precSucc}.
\item \textbf{Trace reflection.}
      Express the invariant as an \code{Fm PTm}; prove quantifier-freeness;
      invoke \code{PRA.qfInd}.
\item \textbf{Reverse demand compiler.}
      For a fixed graph descriptor generated at Lean elaboration time, build
      the concrete \code{PR} reverse fold and prove edge adequacy.
\item \textbf{Node combinators.}
      Implement identity, deterministic map, product, guarded branch,
      bounded search, relational refinement, and sequential composition.
\item \textbf{DAG theorem.}
      Generate the global invariant from the node list and compose the local
      \code{Prf} objects.
\item \textbf{Cubic benchmark.}
      Instantiate the generic trace with rational bisection and export the
      all-precision residual theorem.
\item \textbf{Native transfer.}
      Only after the \PRAAx\ theorem is complete, apply
      \code{pra\_iff\_native} or \code{NFm.ofPRA} to obtain a theorem in the
      native scalar core.
\end{enumerate}

\subsection{Reflection versus theorem generation}

There are two complementary strategies.

\paragraph{Concrete certificate reflection.}
For one closed input and one requested precision, a small verified checker
can reduce the certificate to computation.  A checker soundness theorem then
converts successful evaluation into a proof.

\paragraph{Uniform theorem generation.}
For all inputs, finite computation is not enough.  The generated proof must
contain the base case, step case, and quantifier-free induction.  A
meta-level compiler may construct these proof terms from a graph descriptor,
but the resulting terms are checked by Lean's kernel against the \PRAAx\
rules.

The second strategy is the long-term goal.  It avoids trusting a large tactic:
the compiler may be untrusted because it emits explicit \code{Prf} objects.

\subsection{Formalization status}

The Lean sources supplied here are a detailed blueprint and semantic
prototype.  Lean 4.32.0 was not available in the artifact-building runtime,
so the files were not compiler-checked.  The archive marks every unfinished
object theorem with \code{sorry} and never treats it as established.
The mathematical proofs in this report identify the required invariants and
the exact existing constructors needed to remove those placeholders.
